% typing.tex - Type System and Capability Tracking

\subsection{Type System and Capability Tracking}
\label{sec:typing}

\autoref{fig:typing-rules} presents the full system. %
The type system assigns each term both a type and a capability set. %
The type
describes the value produced; the capability set is a static upper bound on
the authority the term may use. %
The judgment
\[
\typ{\cpdecl}{\tpenv}{t}{T}{\varphi}
\]
means that term $t$ has type $T$ and may access capabilities in $\varphi$. %

\begin{figure}[t]
\centering
\begin{mathpar}
\inferrule*[Right=\rn{T-Const}]
  { }
  {\typ{\cpdecl}{\tpenv}{\tmconst{n}}{\tnat}{\effempty}}

\inferrule*[Right=\rn{T-Var}]
  {\tpenv(x) = T}
  {\typ{\cpdecl}{\tpenv}{x}{T}{\effempty}}

\inferrule*[Right=\rn{T-CVar}]
  {c \in \cpdecl \\
   \cpdecl(c) = T}
  {\typ{\cpdecl}{\tpenv}{\tmcvar{c}}{T}{\eff{c}}}

\\

\inferrule*[Right=\rn{T-Abs}]
  {\typ{\cpdecl}{\tpenv, x{:}T_1}{t}{T_2}{\varphi'}}
  {\typ{\cpdecl}{\tpenv}{\tmabs{x}{T_1}{\varphi}{\varphi}{t}}
       {\tarr{T_1}{T_2}{\varphi}{\varphi}}{\varphi' \effdiff \varphi}}


\inferrule*[Right=\rn{T-App}]
  {\typ{\cpdecl}{\tpenv}{t_1}{\tarr{T_1}{T_2}{\varphi}{\varphi}}{\varphi_1} \\
   \typ{\cpdecl}{\tpenv}{t_2}{T_1}{\varphi_2}}
  {\typ{\cpdecl}{\tpenv}{\tmapp{t_1}{t_2}}{T_2}
       {\varphi \effunion \varphi_1 \effunion \varphi_2}}

\inferrule*[Right=\rn{T-With}]
  {\cpdecl(c) = T_1 \\
   \typ{\cpdecl}{\tpenv}{e_1}{T_1}{\varphi_1} \\
   \typ{\cpdecl}{\tpenv}{e_2}{T_2}{\varphi_2}}
  {\typ{\cpdecl}{\tpenv}{\tmwith{c}{e_1}{e_2}}{T_2}
       {\varphi_1 \effunion (\varphi_2 \effdiff \eff{c})}}
\\

\inferrule*[Right=\rn{T-Allow}]
  {\typ{\cpdecl}{\tpenv}{e}{T}{\varphi'} \\
   \varphi' \effsubset \varphi}
  {\typ{\cpdecl}{\tpenv}{\tmallow{\varphi}{e}}{T}{\varphi'}}

\inferrule*[Right=\rn{T-Sub}]
  {\typ{\cpdecl}{\tpenv}{t}{T}{\varphi} \\
   c \in \cpdecl}
  {\typ{\cpdecl}{\tpenv}{t}{T}{\varphi \effunion \eff{c}}}
\end{mathpar}
\caption{Typing rules for \lambdaCap{}}
\label{fig:typing-rules}
\end{figure}

\rn{T-CVar} is the rule that actually uses a capability: accessing $\tmcvar{c}$
adds exactly $c$ to the term's capabilities. %
\rn{T-Abs} is the central rule. %
If the function body may use capabilities $\varphi'$, and the lambda marks
$\varphi$ as caller-provided, then the lambda expression itself requires only
$\varphi' \effdiff \varphi$ at its definition site. %
In other words, capabilities listed in $\varphi$ are expected from the call site
rather than from closure creation. %
The function type records the caller-provided capability set $\varphi$ above the
arrow. %

\rn{T-App} restores caller responsibility: applying a function of type
$\tarr{T_1}{T_2}{\varphi}{\varphi}$ contributes $\varphi$ to the application's
capabilities. %
This is the static counterpart of the runtime semantics, where the closure does
not capture the capabilities in $\varphi$ and the caller must supply them at
application time. %

\rn{T-With} binds a value to capability $c$ and masks that requirement in the
body, so uses of $c$ inside $e_2$ do not remain obligations of the outer
context. %
\rn{T-Allow} turns dynamic restriction into a static contract: if $e$ checks
with capability set $\varphi'$ and $\varphi' \subseteq \varphi$, then
restricting execution to $\varphi$ cannot remove a capability that the term
needs. %
\rn{T-Sub} is standard weakening. %
